% theme: proving_triangle_similarity
% 叩き台：表・裏とも大問4題。穴埋めは各2題、残りは全文記述。
% 既存図は表示のみ利用し、図ファイルは変更しない。
\MajorQuestionLesson{相似の証明１}
\SetRowSpace{3mm}

\Needspace{75mm}
\TypeQuestion{図で、点 $\mathrm{D},\mathrm{E}$ はそれぞれ辺 $\Lseg{AB},\Lseg{AC}$ 上にあり、$\Lseg{DE}\Parallel\Lseg{BC}$ である。$\Tri{ADE}\Similar\Tri{ABC}$ となることを証明したい。次の証明の空欄をうめなさい。}{triangle_similarity_proof_parallel_fill}
\FigProblemBlock{43mm}{fig/lesson_03_nested}{
\WholeQuestion{%
%「証明」表示を.styでAPI化してください。【】というカッコもつけて
\textbf{証明}\par\smallskip

$\Tri{ADE}$ と $\Tri{ABC}$ において、共通な角であるから、\par
$\Ang{DAE}=$\Blank{ア}\par\smallskip
$\Lseg{DE}\Parallel\Lseg{BC}$ より、平行線の\Blank{イ}は等しいから、\par
$\Ang{ADE}=$\Blank{ウ}\par\smallskip
よって、\Blank{エ}がそれぞれ等しいので、\par
$\Tri{ADE}\Similar\Tri{ABC}$
}{}
}{}

%図を回り込みさせ、文章と図を横並びに。
%穴埋め問題にしてください。
\Needspace{75mm}
\TypeQuestion{}{triangle_similarity_proof_two_sides_angle}
% 図は等角・平行の印なし。長さは文中指定のみ。
\Text{図で、点 $\mathrm{D},\mathrm{E}$ はそれぞれ辺 $\Lseg{AB},\Lseg{AC}$ 上にあり、$\Lseg{AD}=3\mathrm{cm}$、$\Lseg{AE}=4\mathrm{cm}$、$\Lseg{AB}=8\mathrm{cm}$、$\Lseg{AC}=6\mathrm{cm}$ である。$\Tri{ADE}\Similar\Tri{ACB}$ となることを証明しなさい。}
\FigProblemBlock{45mm}{fig/lesson_05_two_sides_angle}{\WholeQuestion{}{}}{}

%もう１題追加

\LessonQuestionBreak

\Needspace{100mm}
\TypeQuestion{図の $\Tri{ABC}$ は $\Lseg{AB}=\Lseg{AC}$ の二等辺三角形である。辺 $\Lseg{AB}$ 上に、$\Ang{BCD}=\Ang{BAC}$ となる点 $\mathrm{D}$ をとった。次の問いに答えなさい。}{triangle_similarity_proof_isosceles_fill}
\FigProblemBlock{35mm}{fig/lesson_04_isosceles}{
\QQ{$\Tri{BCD}\Similar\Tri{BAC}$ となることを示す次の証明の空欄をうめなさい。\par\smallskip
\textbf{証明}\par
$\Tri{BCD}$ と $\Tri{BAC}$ において、仮定より、\par
$\Ang{BCD}=$\Blank{ア}\par\smallskip
点 $\mathrm{D}$ は辺 $\Lseg{AB}$ 上にあるから、\par
$\Ang{CBD}=$\Blank{イ}\par\smallskip
よって、\Blank{ウ}がそれぞれ等しいので、\par
$\Tri{BCD}\Similar\Tri{BAC}$
}{}
}{
%これは穴埋めではありませんが、残していいと思う。
\QQ{前の小問で示した相似と $\Lseg{AB}=\Lseg{AC}$ を用いて、$\Lseg{CD}=\Lseg{BC}$ となることを証明しなさい。}{}
}

\Needspace{110mm}
\TypeQuestion{図の長方形 $\mathrm{ABCD}$ で、辺 $\Lseg{BC}$ 上に点 $\mathrm{F}$ をとり、線分 $\Lseg{AF}$ を折り目として折ると、点 $\mathrm{B}$ は対角線 $\Lseg{AC}$ 上の点 $\mathrm{E}$ に移った。次の問いに答えなさい。}{triangle_similarity_proof_fold_fill}
\FigProblemBlock{42mm}{fig/lesson_04_fold_diagonal}{
\QQ{$\Tri{CEF}\Similar\Tri{CBA}$ となることを示す次の証明の空欄をうめなさい。\par\smallskip
\textbf{証明}\par
$\Tri{CEF}$ と $\Tri{CBA}$ において、折り返しにより、\par
$\Ang{AEF}=$\Blank{ア}$=90^\circ$\par\smallskip
点 $\mathrm{A},\mathrm{E},\mathrm{C}$ は一直線上にあるから、\par
$\Ang{CEF}=90^\circ=\Ang{CBA}$\par\smallskip
また、共通な角であるから、\par
$\Ang{ECF}=$\Blank{イ}\par\smallskip
よって、\Blank{ウ}がそれぞれ等しいので、\par
$\Tri{CEF}\Similar$\Blank{エ}
}{}
}{
%これは穴埋めではありませんが、残していいと思う。
\QQ{前の小問で示した相似を用いて、$\Lseg{CE}\cdot\Lseg{AB}=\Lseg{EF}\cdot\Lseg{BC}$ となることを示しなさい。}{}
}

